\chapter{\Lang{Pendahuluan}{Introduction}}

\section{\Lang{Latar Belakang}{Background}}
\TemplateTodo{\Lang{Jelaskan konteks umum topik, kondisi ideal, masalah nyata, kesenjangan penelitian atau kebutuhan praktis, serta alasan penelitian perlu dilakukan.}{Explain the general context, ideal condition, real problem, research/practical gap, and why the study is necessary.}}

\Lang{Tuliskan latar belakang secara runtut dari gambaran besar menuju masalah khusus \DocTypeNameLower. Gunakan sumber ilmiah yang dapat diverifikasi, misalnya artikel jurnal, prosiding, buku, standar, atau dokumen resmi \cite{exampleReference}.}{Write the background coherently from the broader context toward the specific \DocTypeNameLower\ problem. Use verifiable scholarly sources such as journal articles, conference papers, books, standards, or official documents \cite{exampleReference}.}

\TemplateTip{\Lang{Hindari langsung membahas detail teknis tanpa menjelaskan masalah. Setiap klaim penting sebaiknya didukung sitasi.}{Avoid jumping into technical details before explaining the problem. Support important claims with citations.}}

\section{\Lang{Rumusan Masalah}{Problem Statement}}
\TemplateTodo{\Lang{Ubah masalah penelitian menjadi pertanyaan yang spesifik, terukur, dan dapat dijawab melalui metodologi.}{Convert the research problem into specific, measurable questions that can be answered through the methodology.}}

\begin{enumerate}
  \item \Lang{Bagaimana merancang pendekatan untuk menyelesaikan masalah utama pada objek penelitian?}{How can an approach be designed to address the main problem in the research object?}
  \item \Lang{Bagaimana mengevaluasi kinerja pendekatan yang diusulkan berdasarkan metrik yang sesuai?}{How can the proposed approach be evaluated using appropriate metrics?}
  \item \Lang{Faktor apa saja yang memengaruhi hasil penelitian berdasarkan analisis yang dilakukan?}{What factors influence the research results based on the analysis?}
\end{enumerate}

\section{\Lang{Batasan Masalah}{Scope and Limitations}}
\TemplateTodo{\Lang{Tuliskan ruang lingkup agar penelitian tidak melebar. Batasan dapat berupa data, metode, lokasi, periode, variabel, alat, atau asumsi.}{Define the scope so the study remains focused. The scope may cover data, method, location, period, variables, tools, or assumptions.}}

\section{\Lang{Tujuan Penelitian}{Research Objectives}}
\TemplateTodo{\Lang{Selaraskan tujuan dengan rumusan masalah. Gunakan kata kerja yang dapat diverifikasi seperti menganalisis, merancang, mengimplementasikan, mengevaluasi, atau membandingkan.}{Align objectives with the problem statements. Use verifiable verbs such as analyze, design, implement, evaluate, or compare.}}

\section{\Lang{Manfaat Penelitian}{Research Benefits}}
\TemplateTodo{\Lang{Jelaskan manfaat teoretis dan praktis. Hindari klaim terlalu besar; tuliskan kontribusi yang realistis sesuai skala \DocTypeNameLower.}{Explain theoretical and practical benefits. Avoid exaggerated claims; state realistic contributions for the scale of the \DocTypeNameLower.}}